% Part: sets-functions-relations
% Chapter: arithmetization
% Section: reflections

\documentclass[../../../include/open-logic-section]{subfiles}

\begin{document}

\olfileid{sfr}{arith}{ref}
\olsection{কয়েকটি দার্শনিক প্রতিফলন}

কারিগরি খুঁটিনাটির কথা এতটুকুই। কিন্তু এগুলি কী অর্জন করল?

বেশ অবিসংবাদিতভাবেই বলা যায়, এগুলি আমাদের কিছু {সুন্দর} বিশুদ্ধ গণিত দিয়েছে। উপরন্তু, এতে কয়েকটি গভীর ধারণাগত অর্জন ছিল। পূর্ণতা ধর্মই যে বাস্তব ও মূলদ সংখ্যার মধ্যে গুরুত্বপূর্ণ পার্থক্য প্রকাশ করে, তা দেখতে পাওয়া ছিল এক গভীর অন্তর্দৃষ্টি। আবার, ডেডেকিন্ড কর্তন হিসেবে বাস্তব সংখ্যার স্পষ্ট নির্মাণ গাণিতিক বিশ্লেষণের বিষয়বস্তুকে দৃঢ় ভিত্তি দেয়। \emph{পূর্ণ ক্রমিত ক্ষেত্র}-এর ধারণা যে সঙ্গতিপূর্ণ তা আমরা জানি, কারণ কর্তনগুলি ঠিক এমন একটি ক্ষেত্র গঠন করে।

এত কিছু সত্ত্বেও, এই অর্জনগুলি সম্পর্কে কয়েকটি সংশয় প্রকাশ করা উচিত।

প্রথমত, বাস্তব সংখ্যাকে কর্তনের মাধ্যমে ভাবা তাদের পরিচিত (সম্ভবত অসীম) দশমিক বিস্তারের মাধ্যমে ভাবার চেয়ে \emph{বেশি} কঠোর কি না, তা স্পষ্ট নয়। বাস্তব সংখ্যার এই শেষোক্ত ``নির্মাণ'' কোশি অনুক্রমের সাহায্যে বাস্তব সংখ্যা নির্মাণের সঙ্গে কিছুটা সাদৃশ্যপূর্ণ; কিন্তু আসলে সপ্তদশ শতাব্দীর গোড়া থেকেই গণিতবিদেরা এর মূল কথা জানতেন (\olref[cauchy]{sec} দেখো)। কঠোরতার প্রকৃত বৃদ্ধি এসেছিল বাস্তব সংখ্যার পূর্ণতা ধর্ম আছে—এই উপলব্ধি থেকে; নির্দিষ্ট কিছু সেট হিসেবে বাস্তব সংখ্যা নির্মাণের সামর্থ্যটি হয়তো নিজে থেকে খুব আকর্ষণীয় নয়।

পূর্ণসংখ্যা বা মূলদ সংখ্যার (অনেক সহজ) সেটতাত্ত্বিক সংখ্যা-নির্মাণ ওই ক্ষেত্রগুলিতে কঠোরতা বাড়ায় কি না, তা আরও অস্পষ্ট। এখানে একটি সহজ পর্যবেক্ষণ করা দরকার। স্বাভাবিক সংখ্যার ক্রমযুগলের তুল্যতা শ্রেণি হিসেবে পূর্ণসংখ্যা \emph{নির্মাণ} করার পরে, পূর্ণসংখ্যার ক্রমযুগলের তুল্যতা শ্রেণি হিসেবে মূলদ সংখ্যা নির্মাণ করার পরে, এবং মূলদ সংখ্যার সেট হিসেবে বাস্তব সংখ্যা নির্মাণ করার পরে, আমরা নির্মাণগুলির কথা সঙ্গে সঙ্গে \emph{ভুলে যাই}। বিশেষত, কোনো গাণিতিক প্রমাণে কেউ কখনো এই নির্মাণগুলি \emph{প্রয়োগ} করতে চাইবে না (অবশ্য নির্মাণগুলি যেমন আচরণ করার কথা তেমনই করে—এটি প্রমাণ করার সময় ছাড়া)। স্বাভাবিক সংখ্যার সেটের সেটের সেটের সেটের সেটের সেটের সেটের কোনো সেট নিয়ে কথা বলার চেয়ে সরাসরি একটি বাস্তব সংখ্যা নিয়ে কথা বলা অনেক সহজ।
%(কারণ বাস্তব সংখ্যা মূলদ সংখ্যার সেট হিসেবে নির্মিত; আর মূলদ সংখ্যা পূর্ণসংখ্যার ক্রমযুগলের তুল্যতা শ্রেণি হিসেবে, অর্থাৎ পূর্ণসংখ্যার সেটের সেটের সেট হিসেবে নির্মিত; ইত্যাদি।)
%2/3 = [2,3]
%	অর্থাৎ {<2,3>, ... }
%	{ {{2},{2,3}}, ... }
%
%বাস্তব সংখ্যা হলো
%	মূলদ সংখ্যার সেট
%
%মূলদ সংখ্যা হলো
%	পূর্ণসংখ্যার সেটের সেটের সেট
%
%পূর্ণসংখ্যা হলো
%	স্বাভাবিক সংখ্যার সেটের সেটের সেট

এই সংজ্ঞাগুলি পূর্ণসংখ্যা, মূলদ সংখ্যা বা বাস্তব সংখ্যা \emph{অধিবিদ্যাগতভাবে আসলে কী}, তা বলে—এটি সবচেয়ে বেশি সন্দেহজনক। অর্থাৎ, বাস্তব সংখ্যাগুলি (ধরা যাক) নির্দিষ্ট কিছু সেট (সেটের সেটের সেট\ldots) \emph{হয়}—এ দাবি সন্দেহজনক। এমন মতের প্রধান বাধা হলো, নির্মাণটি বহু ভিন্নভাবে করা যেত। বাস্তব সংখ্যার ক্ষেত্রে সত্যিই আকর্ষণীয় কয়েকটি ভিন্ন নির্মাণ আছে (\olref[cauchy]{sec} দেখো)। তবে ভিন্ন নির্মাণ পাওয়ার একেবারেই তুচ্ছ একটি উপায় এখানে: \olref[sfr][rel][ref]{sec}-এর মতো আমরা ক্রমযুগলকে একটু অন্যভাবে সংজ্ঞায়িত করতে পারতাম; ক্রমযুগলের সেই বিকল্প ধারণা ব্যবহার করলে আমাদের নির্মাণগুলি আগের মতোই নিখুঁতভাবে কাজ করত, কিন্তু শেষে আমরা ভিন্ন বস্তু পেতাম। সেই কারণে পূর্ণসংখ্যা, মূলদ সংখ্যা ও বাস্তব সংখ্যার পরস্পর-প্রতিদ্বন্দ্বী অনেক সেটতাত্ত্বিক নির্মাণ আছে। আর তখন পূর্ণসংখ্যা (ধরা যাক) \emph{ওই} সেটগুলি না হয়ে \emph{এই} সেটগুলিই—এ দাবি করা নিতান্ত খামখেয়ালি (এবং বিব্রতকর) হবে। (\olref[sfr][rel][ref]{sec}-এর মতো, এটি \citealt{Benacerraf1965}-এর সুবিখ্যাত একটি যুক্তির নিদর্শন।)

আরও একটি কথা তোলা দরকার: আমাদের নির্মাণগুলির মধ্যে বেশ \emph{অদ্ভুত} কিছু আছে। আমরা স্বাভাবিক সংখ্যা দিয়ে শুরু করেছি। তারপর পূর্ণসংখ্যা নির্মাণ করি, এবং ``পূর্ণসংখ্যার $0$'', অর্থাৎ $ \equivrep{0,0}{\Intequiv}$ নির্মাণ করি। কিন্তু $0 \neq
\equivrep{0,0}{\Intequiv}$। বস্তুত, আমাদের নির্মাণ অনুসারে কোনো স্বাভাবিক সংখ্যাই পূর্ণসংখ্যা \emph{নয়}। কিন্তু এটি অত্যন্ত স্বজ্ঞাবিরোধী মনে হয়। আসলে \olref[sfr][set][imp]{sec}-এ বিশেষ যুক্তি ছাড়াই আমরা দাবি করেছিলাম যে $\Nat \subseteq \Rat$। নির্মাণগুলি যদি সংখ্যা ঠিক \emph{কী} তা বলে, তবে ওই দাবি তুচ্ছভাবেই মিথ্যা ছিল।

তাহলে একটু দূরে দাঁড়িয়ে দেখলে আমরা কোথায় পৌঁছাই? অনানুষ্ঠানিক সেটতত্ত্বে কাজ করে এবং স্বাভাবিক সংখ্যাগুলিকে ধরে নিয়ে আমরা পূর্ণসংখ্যা, মূলদ সংখ্যা ও বাস্তব সংখ্যাকে নির্দিষ্ট কিছু সেট হিসেবে \emph{ধরতে} পারি। সেই অর্থে এই সত্তাগুলির তত্ত্বগুলিকে একটি সেটতত্ত্বের মধ্যে \emph{অন্তর্ভুক্ত} করতে পারি। কিন্তু এই অন্তর্ভুক্তির দার্শনিক তাৎপর্য মোটেই সরল নয়।

অবশ্যই, এটাই শেষ কথা নয়!{} বক্তব্যটি শুধু এই: বাস্তব সংখ্যার সেটতাত্ত্বিক সংখ্যা-নির্মাণের \emph{গভীর} দার্শনিক তাৎপর্য \emph{আছে} দেখাতে আরও কিছু \emph{দার্শনিক} যুক্তি লাগবে।

%আরও একটি কথা: এই পুরো অধ্যায়ে আমরা স্বাভাবিক সংখ্যাগুলিকে স্রেফ আমাদের কাছে \emph{দেওয়া} আছে বলে ধরেছি। কিন্তু তাদের নিয়ে আমরা কী করতে পারি? সেটিই পরের অধ্যায়ের বিষয়।

\end{document}
